% 100_class_theory.tex -- Class Theory Overview & Morse-Kelley Comprehension / Separation.
% Tickets 7335 & 7454: Rigorous Morse-Kelley class comprehension schema, sethood
% distinction, and class separation principles (audited by Claude Opus 4.6).
\def\topic{Class Theory Overview}
\ifdefined\MyHeaderIncluded
\else
\def\qwertyflag{true}

\ifdefined\MyHeaderIncluded

\else
  \newcommand{\foobar}{FooBar}
  \def\MyHeaderIncluded{true}
% if is till the end of the file

\documentclass[10pt,a4paper]{article}
\usepackage[utf8]{inputenc}
%\usepackage[german]{babel}
%\usepackage[cyrillic]{babel}
\usepackage[T1]{fontenc}
%\usepackage[T2A]{fontenc}
%\usepackage[english, russian]{babel}
\usepackage{amsmath}
\usepackage{amsfonts}
\usepackage{amssymb}
\usepackage{amsthm}
\usepackage{mathtools} %coloneqq

\usepackage{bussproofs} %

%\renewcommand{\thesubsection}{\arabic{subsection}}
%\newtheorem{thm}{Theorem}[section]
%\renewcommand{\thethm}{\arabic{subsection}.\arabic{thm}}

\usepackage{thmtools}
\usepackage{hyperref}
\hypersetup{
    colorlinks=true,
    linkcolor=blue,
    filecolor=magenta,      
    urlcolor=cyan,
    pdftitle={Overleaf Example},
    %pdfpagemode=FullScreen,
    }
%\usepackage{stix-type1}%\turnediota
\theoremstyle{definition}

\newtheorem{constructor}{Constructor}[section]
\newtheorem{iaxiom}{Axiom}
\newtheorem{axiom}{Axiom}
\newtheorem*{axiom*}{Axiom}
\newtheorem{definition}{Definition}[section]
\newtheorem*{definition*}{Definition}
\newtheorem{remark}{Remark}[section]
\newtheorem*{remark*}{Remark}
\newtheorem{recall}{Recall}[section]
\newtheorem*{recall*}{Recall}

\newtheorem{lemma}{Lemma}[section]
\newtheorem{theorem}{Theorem}[section]
\newtheorem*{theorem*}{Theorem}
\newtheorem{note}{Note}[section]
\usepackage{etoolbox}

%\newcommand*{\@Dashv}[2]{%
%  \reflectbox{$\m@th#1#2$}%
%}
\newcommand*{\separation}[3]{% 1)binder, 2)set, 3)condition
  \{#1\in #2: #3\}
}

% https://www.freecodecamp.org/news/simple-regex-tricks-for-beginners-3acb3fa257cb/
%  \\mathcal\{T\} |---> \topology
%  \\mathcal\{([^\}]*)\}   |--->  $1
\newcommand{\AND}{\ \ \&\ \ }
\newcommand{\any}{\mathrm{\ any}}
\newcommand{\borel}{\mathcal{B}}
\newcommand{\can}{\mathrm{can}}
\newcommand{\canonize}{\mathrm{canonize}}
\newcommand{\closure}{\mbox{cl}}
\newcommand{\complex}{\mathbb{C}}
\newcommand{\converges}{\longrightarrow}
\newcommand{\countable}{\mbox{\ countable}}
\newcommand{\ldefined}{\mbox{\ defined}}
\newcommand{\dom}{\mbox{dom}}
\newcommand{\Dom}{\mbox{Dom}}
\newcommand{\Done}{\noindent \textbf{Done.}\\}
\newcommand{\lexists}{\mbox{\ exists}}
\newcommand{\lforall}{\mbox{\ forall}}
\newcommand{\lunique}{\mbox{\ unique}}
\newcommand{\family}{\mathrm{\ family}}
\newcommand{\Fm}{\mathrm{Fm}}
\newcommand{\FuncSymb}{\mathrm{FuncSymb}}
\newcommand{\function}{\mathrm{\ function}}
\newcommand{\functional}{\mbox{\ functional}}
\newcommand{\IFF}{\ \Longleftrightarrow\ }
\renewcommand{\iff}{\ \longleftrightarrow\ }
\newcommand{\IMPLIES}{\ \Longrightarrow\ }
\renewcommand{\implies}{\ \longrightarrow\ }
\newcommand{\Inf}{\mathrm{Inf}}
\newcommand{\injective}{\mbox{\ injective}}
\newcommand{\injection}{\mbox{\ injection}}
\newcommand{\interior}{\mbox{int}}
\newcommand{\intersection}{\bigcap}
\newcommand{\isempty}{\mathrm{\ empty}}
\newcommand{\LB}{\mathrm{LB}}
\newcommand{\nghbset}{\mathcal{N}}
\newcommand{\nonempty}{\mathrm{\ nonempty}}
\newcommand{\Ord}{\mathrm{Ord}}
\newcommand{\power}{\mathcal{P}}
\newcommand{\pr}{\mathrm{pr}}
\newcommand{\fst}{\mathrm{fst}}
\newcommand{\snd}{\mathrm{snd}}
\newcommand{\PredSymb}{\mathrm{PredSymb}}
\newcommand{\Proof}{\noindent \textbf{Proof.}\\}
\newcommand{\ran}{\mathrm{ran}}
\newcommand{\Ran}{\mathrm{Ran}}
\newcommand{\real}{\mathbb{R}}
\newcommand{\relation}{\mathrm{\ relation}}
\newcommand{\set}{\mathrm{\ set}}
\newcommand{\singleton}{\mathrm{\ singleton}}
\newcommand{\semantics}{\mathrm{semantics}}
\newcommand{\semialgebra}{\mathcal{S}}
\newcommand{\Sup}{\mathrm{Sup}}
\newcommand{\surjective}{\mbox{\ surjective\ on\ }}
\newcommand{\Term}{\mathrm{Term}}
\newcommand{\Type}{\mathrm{Type}}
\newcommand{\topology}{\mathcal{T}}
\newcommand{\total}{\mbox{\ total\ on\ }}
\newcommand{\transitive}{\mathrm{\ transitive}}
\newcommand{\union}{\bigcup}
\newcommand{\lif}{\mathrm{if}}
\newcommand{\lthen}{\mathrm{then}}
\newcommand{\lelse}{\mathrm{else}}
\newcommand{\lexception}{\mathrm{exception}}
\newcommand{\lundef}{\mathrm{undef}}
\newcommand{\univocal}{\mathrm{\ univocal}}
\newcommand{\univocalat}{\mathrm{\ univocal\ at\ }}
\newcommand{\UB}{\mathrm{UB}}
\newcommand{\Var}{\mathrm{Var}}

\newcommand{\notvarphi}{\phi}
\renewcommand{\phi}{\varphi}

\newcommand{\comment}[1]{\quad\mbox{(#1)}}
% local commands
%\preto\section{\undef\blipo} %etoolbox
\preto\subsection{\undef\blipo} %etoolbox
% https://tex.stackexchange.com/questions/105684/limit-macro-newcommand-scope-to-section
\usepackage{graphicx}
\usepackage{authblk}

\setlength\parindent{0pt} % whole doc of \noindent

%\usepackage{mathgl}

\fi % guard
\ifdefined\topic
\else
\def\topic{No Titel}
\fi
\title{\topic}
\author[1]{Georgy Dunaev}
\begin{document}
\maketitle
\fi
\section{Class Theory Overview}
\label{sec:class_theory_overview}

This section summarizes the Morse--Kelley (MK) and
von Neumann--Bernays--G\"odel (NBG) principles used in the classical
class-theoretic chapters. It connects the axiomatic kernel of
Section~\ref{sec:axioms} with family evaluation, sethood, and the
union and intersection operations.

\subsection{Classes, Sets, and Proper Classes}
\label{subsec:classes_sets_proper}

Every mathematical object in the first-order discourse is a \emph{class}. The universal class of all sets is defined by
\[
V \eqdef \cc{x}{\top}.
\]
A class $x$ is a \emph{set} if and only if it belongs to some class ($x \in V \iff \exists Y,\ x \in Y$). Every set is a class; a \emph{proper class} is a class $X$ that is not a set ($X \notin V$). Consequently, only sets can appear on the left-hand side of membership ($z \in X \implies z \in V$).

\subsection{Morse--Kelley Class Comprehension and Separation}
\label{subsec:mk_comprehension_separation}

For an admitted formula $\phi(x)$, the class comprehension term $\cc{x}{\phi(x)}$ satisfies the comprehension equivalence:
\[
\forall z,\quad z \in \cc{x}{\phi(x)} \iff (z \in V) \land \phi(z).
\]
The default sort of variables is $\mathsf{Class}$. Full MK permits class
quantification inside $\phi$; NBG comprehension restricts its internal
quantifiers to sets while allowing class parameters. In one-sorted syntax,
a set quantifier can be written as a class quantifier restricted to $V$.

Notice the mandatory sethood guard $z \in V$: substituting the Russell class $R \eqdef \cc{x}{x \notin x}$ yields $R \in R \iff (R \in V) \land (R \notin R)$, which proves $\lnot(R \in V)$ (i.e., $R$ is a proper class) rather than a contradiction.

Given any set $a \in V$ and any class $P$ (or formula $\phi$), the \emph{Separation Principle} (Subsets / Intersection with a Set) states that the intersection
\[
a \cap P = \cc{x \in a}{x \in P} = \cc{x}{(x \in a) \land (x \in P)}
\]
is again a set: $(a \in V) \implies (a \cap P \in V)$. In the axiom presentation used here, this is the heredity principle.
In presentations based on limitation of size it can instead be derived
from the relevant size axioms. A derivation of limitation of size must keep
track of which of these principles it assumes, to avoid circularity.

\subsection{What the minimal comprehension principles do not imply}
\begin{theorem}[A finite model of the minimal relational fragment]
\label{thm:minimal_class_countermodel}
Consider classical one-sorted first-order logic in the language of equality
and membership, with sethood defined by
$\operatorname{Set}(x)\iff\exists Y\,(x\in Y)$. Take the following
nonlogical assumptions only: class extensionality, the schema of
set-guarded class comprehension, and existence of at least one set.
This fragment has a model with two class objects and only one set.
Consequently it does not imply the usual pairing or infinity axioms.
\end{theorem}
\Proof
Let the domain consist of two distinct objects $e,v$, and interpret
membership by the single ordered pair $(e,v)$. Thus $e$ has no members,
$v$ has the sole member $e$, and the objects satisfying $\operatorname{Set}$
are exactly $e$. Extensionality holds because the two objects have
different members; the existence-of-a-set assertion holds because
$e\in v$.

Fix any comprehension formula $\phi(x)$ and any assignment of its
parameters. The required extension
$\{x:\operatorname{Set}(x)\land\phi(x)\}$ is either empty or $\{e\}$,
according as $\phi(e)$ is false or true in this structure. It is
represented respectively by $e$ or $v$. This verifies every instance
of the schema, including formulas with unrestricted quantifiers over
the two class objects. The comprehension output variable is, as usual,
fresh for $\phi$; no satisfaction predicate is added to the language.

The only set is the empty object $e$. A set containing $e$, as required
by pairing applied to $e,e$, would have to be $v$, which is not a set.
Thus pairing fails. No set contains even the empty set, so no inductive
set exists and the usual infinity axiom fails. Soundness of first-order
logic gives the nonderivability conclusions.
\Done

This is a diagnostic for exactly the listed relational assumptions, not
a model of every system proposed in this book. Stronger principles about
constructors, typed interpretation, reflection, or additional set existence
can exclude this model. Any proposed interpretation of ZFC in a richer
system must identify those additional principles and prove the required
interpretation properties. Merely naming unrestricted class comprehension
and one set does not supply the missing set-existence axioms.

\begin{theorem}[Finite closure principles do not supply infinity]
Work externally in ZF set theory, with its axiom of infinity. There is a
set-sized model of the one-sorted class language satisfying full
set-guarded class comprehension, extensionality, empty set, pairing,
union, powerset, separation, and replacement for class functions, but
not the axiom of infinity. In this model every internal set is externally
finite. The result concerns these explicitly listed principles, not an
unspecified stronger reflection or universe principle.
\end{theorem}
\Proof
Let $H=V_\omega$ be the set of hereditarily finite sets in the external
theory, and take the class-object domain to be $\mathcal P(H)$.
Interpret membership as ordinary membership restricted to that domain.
Since $H$ is transitive, every member of $H$ is a subset of $H$ and
therefore an object of the structure. An object $x$ belongs to another
object exactly when $x\in H$: one implication follows because every
object is a subset of $H$; for the other, the object $\{x\}$ witnesses
membership. Thus internal sethood is precisely membership in $H$.

All members of an object are themselves in the domain, so external
extensionality gives internal class extensionality. For any formula
$\phi(x)$ and parameters, including formulas quantifying over all class
objects, its required comprehension extension is
\[
\{x\in H:\phi(x)\text{ holds in the structure}\}.
\]
This is a subset of $H$, hence an object. The verification is external
and applies instance by instance; it introduces no internal truth
predicate.

The empty set belongs to $H$. Pairs of members of $H$, unions of
members of $H$, and powersets of members of $H$ again belong to $H$:
all these constructions yield finite sets of hereditarily finite
objects. Every subclass of an internal set $a$ is an external subset
of a finite hereditarily finite set, hence is again in $H$. Finally,
a class function with internal set domain $a$ has only finitely many
inputs. Its image is a finite set of members of $H$, and therefore is
in $H$. This proves class replacement, regardless of how the class
function was defined. Ordinary finite graphs are also in $H$.

If an internal set were inductive, it would contain the external empty
set and, by external induction, every finite von Neumann ordinal.
Those ordinals are distinct, contradicting external finiteness of the
internal set. Thus infinity fails.
\Done

\begin{corollary}[Even class limitation of size does not force infinity]
The preceding hereditarily finite model also satisfies the following
class-size principle: an internal class is proper if and only if it is
in internal bijection with the universal class. Thus adding this
principle to the listed closure axioms still does not imply infinity.
\end{corollary}
\Proof
Externally, $H$ has a fixed bijection with $\omega$. One explicit coding
is the recursive finite-set code
$c(x)=\sum_{y\in x}2^{c(y)}$. Induction on finite ranks and uniqueness
of binary expansion prove injectivity. Conversely, strong induction on
$n$ constructs a set with code $n$ from the sets coding its binary
exponents, all of which are smaller than $n$ when $n>0$. Hence $c$ is
bijective. This coding and the increasing enumeration below use no
external choice axiom. The natural numbers here are external indices;
they are not being asserted to form an internal set.

An object $A\subseteq H$ is an internal set exactly when it is
externally finite: every finite subset of $H$ is hereditarily finite,
and every member of $H$ is finite. If $A$ is proper, its codes form an
infinite subset of $\omega$. Listing them in increasing order gives an
external bijection $H\to A$. Its graph consists of ordinary ordered
pairs of hereditarily finite sets, which are themselves hereditarily
finite; hence that graph is a subset of $H$ and is an internal class
object. It witnesses the required internal bijection. Conversely, an
internal bijection from $H$ to an internal set would externally biject
an infinite set with a finite one, which is impossible.

This also verifies the book's surjective-image formulation of limitation
of size: for proper $A$, the inverse of the coded bijection $H\to A$
is an internal class function $T$ with $H=T[A]$. If $A$ is an
internal set, every set-valued class function has finite image on $A$,
so no such image can contain the infinite universal class $H$.
\Done

\begin{corollary}[Regularity and global choice in the same model]
The hereditarily finite model also satisfies class foundation and has a
class choice function on all nonempty internal sets. These principles
therefore do not repair its failure of infinity.
\end{corollary}
\Proof
For a nonempty internal class $A\subseteq H$, the external finite ranks
of its members form a nonempty subset of $\omega$ and have a least
value. Choose one member $a$ of that rank. If $b\in a\cap A$, then
$b$ has strictly smaller rank, a contradiction. Thus $a\cap A$ is
empty. This verifies class foundation, and in particular set regularity.
It makes no assertion that the class of all finite ordinals is an
internal set.

For a nonempty $a\in H$, select its member with least Ackermann code
$c(x)$. The code is injective, so the selected member is unique.
The graph of this operation consists of ordered pairs of members of
$H$, hence is a subset of $H$ and an internal class. Its domain is
exactly the nonempty internal sets and its value at $a$ belongs to $a$.
This is the claimed global choice function. All choices used to define
it were uniquely determined by the fixed coding.
\Done

The two models isolate different missing assumptions. The two-object
model shows that comprehension and one set do not even give pairing.
The hereditarily finite model shows that adding the listed set-closure
principles still does not give infinity. These are relative model
constructions in an external theory; neither is a claim that the
book's stronger proposed systems have the same models. They provide
concrete tests for any claimed derivation of set-existence principles.

\subsection{Ordered Pairs of Classes: Dunaev Tagged Ordered Pairs vs.\ Kuratowski Pairs}
\label{subsec:dunaev_vs_kuratowski_pairs}

In Morse-Kelley class theory, class comprehension ranges over sets: a class
$X$ is a member of $\cc{x}{\phi(x)}$ only when $X \in V$. Consequently, the
singleton $\{A\} = \cc{x}{x = A}$ equals the empty class $\emptyclass$ whenever
$A \notin V$ is a proper class, and the Kuratowski pair
$\kuratowskipair{A}{B} = \{\{A\}, \{A, B\}\}$ collapses to $\{\emptyclass\}$
whenever $A$ and $B$ are proper classes, destroying injectivity of projections.
Therefore, Kuratowski pairs must never be used for proper classes or general
class-level pairing. Instead, class-level ordered pairs use Dunaev Tagged
Ordered Pairs (Sections~\ref{sec:ordered_pairs_foundations} and~\ref{sec:dunaev_tagged_ordered_pair}):
\[
\pair{A}{B} \eqdef \cc{z}{\exists x,\ (z = \kuratowskipair{\emptyclass}{x} \land x \in A) \lor (z = \kuratowskipair{\{\emptyclass\}}{x} \land x \in B)}.
\]
Because the tag coordinates $\emptyclass$ and $\{\emptyclass\}$ and the element
witnesses $x \in A \subseteq V$ and $x \in B \subseteq V$ are always sets,
every member of $\pair{A}{B}$ is a set in $V$ regardless of whether $A$ and $B$
are sets or proper classes, and the projections
$\Fst(P) = \cc{x}{\kuratowskipair{\emptyclass}{x} \in P}$ and
$\Snd(P) = \cc{x}{\kuratowskipair{\{\emptyclass\}}{x} \in P}$ recover $A$ and $B$
uniquely for all classes $A, B$.

\subsection{Binders: Quantifiers, Class Comprehension, and Descriptors}
\label{subsec:binders_overview}

A \emph{binder} is a syntactic construction that binds a variable.
The following list is a guide to the notation encountered in different
sections, not a claim that every listed binder has the same semantics in
all of them. The default class sort applies to the classical setting;
partial and uncertain semantics supply their own interpretation rules.
\begin{enumerate}
  \item \textbf{Formula-Valued Logical Quantifiers}: universal ($\forall x,\ \phi(x)$), existential ($\exists x,\ \phi(x)$), uniqueness ($\exists! x,\ \phi(x)$ and $!x,\ \phi(x)$), and bounded quantifiers ($\forall (x \in Y),\ \phi(x)$ and $\exists (x \in Y),\ \phi(x)$).
  \item \textbf{Class Comprehension Binder ($\kappa$-binder)}: $\cc{x}{\phi(x)}$ (and generalized comprehension $\ccc{\vec{x}}{t(\vec{x})}{\phi(\vec{x})}$), binding a class variable $x:\mathsf{Class}$ to form a class term.
  \item \textbf{Definite Description ($\iota$-binder)}: $\iota x,\ \phi(x)$, denoting the unique class satisfying $\phi(x)$ when $\exists! x,\ \phi(x)$ holds (with its behavior in other cases specified by the local descriptor convention).
  \item \textbf{Choice and Partial Descriptors} ($\tau$, $\epsilon$,
  $\varepsilon$; Section~\ref{sec:descriptors}): the descriptor experiment
  distinguishes an arbitrary-class default for $\tau$, an empty-class
  default for $\epsilon$, and no default for partial $\varepsilon$.
  These are local conventions. In particular, an undefined expression is
  not automatically identified with a term denoting the empty class.
\end{enumerate}

\subsection{Finite Class-Lists over Dunaev Tagged Ordered Pairs (\texttt{3.51})}
\label{subsec:class_lists_dunaev}

Catalog Item \texttt{3.51} (\texttt{DefClassLists}) defines the class $\mathrm{Lists}(A)$ of finite lists with entries in a class $A$ by impredicative intersection over all classes $C$ containing the empty sentinel $\emptyclass$ and closed under prepending elements of $A$ via Dunaev Tagged Ordered Pairs:
\[
\mathrm{Lists}(A) \eqdef \bigcap \Big\{ C \subseteq \V \;\Big|\; \emptyclass \in C \land \forall h \in A\,\forall t \in C,\ \pair{h}{t} \in C \Big\}.
\]
Because $\pair{h}{t}$ is a set whenever $h, t \in \V$ and projections $\Fst(\pair{h}{t}) = h$ and $\Snd(\pair{h}{t}) = t$ are unconditionally injective, $\mathrm{Lists}(A)$ satisfies the Peano-style list constructors and structural induction principle:
\begin{enumerate}
  \item \textbf{Constructor Closure and Disjointness}: $\emptyclass \in \mathrm{Lists}(A)$, and for all $h \in A$ and $t \in \mathrm{Lists}(A)$, $\pair{h}{t} \in \mathrm{Lists}(A)$ with $\pair{h}{t} \neq \emptyclass$.
  \item \textbf{Structural Induction and Transport}: For any class property $\Phi(u)$, if $\Phi(\emptyclass)$ holds and $\forall h \in A\,\forall t \in \mathrm{Lists}(A),\ (\Phi(t) \implies \Phi(\pair{h}{t}))$, then $\forall u \in \mathrm{Lists}(A),\ \Phi(u)$.
\end{enumerate}

\subsection{Indexed-Union Closure of Tight Grothendieck Universes (\texttt{3.68}, \texttt{3.69})}
\label{subsec:tight_universe_indexed_union}

In Catalog Item \texttt{3.68} (\texttt{DefGrothendieckUniverse}), a class $\mathcal{U}$ is defined as a \emph{tight Grothendieck universe} when it is transitive ($\forall x \in \mathcal{U},\ x \subseteq \mathcal{U}$), closed under powersets ($\forall x \in \mathcal{U},\ \mathcal{P}(x) \in \mathcal{U}$), and subset-tight ($\forall x \in \mathcal{U}\,\forall y \subseteq \mathcal{U},\ |y| \le |x| \implies y \in \mathcal{U}$).

\begin{proposition}[Indexed-Union Closure and Universal Class Universe Theorem (\texttt{3.68}, \texttt{3.69})]
\label{prop:tight_universe_indexed_union}
Every tight Grothendieck universe $\mathcal{U}$ is closed under set-indexed unions: for every index set $I \in \mathcal{U}$ and every family $A$ assigning to each $i \in I$ a set $A(i) \in \mathcal{U}$, the union $\bigcup_{i \in I} A(i)$ belongs to $\mathcal{U}$. Moreover, the universal class $\V$ satisfies \texttt{3.68} (\texttt{3.69 ThmUniversalClassIsUniverse}).
\end{proposition}
\begin{proof}
Since $I \in \mathcal{U}$ and $A(i) \in \mathcal{U}$ for each $i \in I$, transitivity of $\mathcal{U}$ implies $A(i) \subseteq \mathcal{U}$ for every $i \in I$, and hence the range $R = \cc{y}{\exists i \in I,\ y = A(i)}$ is a subclass of $\mathcal{U}$ with $|R| \le |I|$. By subset-tightness of $\mathcal{U}$ applied to $I \in \mathcal{U}$ and $R \subseteq \mathcal{U}$, we obtain $R \in \mathcal{U}$. Applying single-set union closure (derived from powerset closure and subset-tightness inside $\mathcal{U}$) yields $\bigcup R = \bigcup_{i \in I} A(i) \in \mathcal{U}$. For $\mathcal{U} = \V$ (Item \texttt{3.69}), transitivity holds because elements of sets are sets, powerset closure holds by the Powerset Axiom, and subset-tightness and indexed-union closure hold by the Axiom of Limitation of Size (Replacement) and the Union Axiom.
\end{proof}

\subsection{Regular-Open Completion and Boolean-Valued Models over Dunaev-Tagged Names}
\label{subsec:ro_completion_boolean_valued_models}

Given a forcing preorder $(\mathbb{P}, \le, 1_{\mathbb{P}})$, for any subset $U \subseteq \mathbb{P}$ let $U^\perp \eqdef \cc{q \in \mathbb{P}}{\forall p \in U,\ \neg(p \parallel q)}$ denote the incompatible complement, and let $U^{\perp\perp} \eqdef (U^\perp)^\perp$ denote the regular-open closure. The regular-open completion $\mathrm{ro}(\mathbb{P}) \eqdef \cc{U \subseteq \mathbb{P}}{U = U^{\perp\perp}}$ ordered by inclusion forms a complete Boolean algebra $(\mathbb{B}, \wedge, \vee, \neg, 0_{\mathbb{B}}, 1_{\mathbb{B}})$ with $U \wedge W = U \cap W$, $U \vee W = (U \cup W)^{\perp\perp}$, and $\neg U = U^\perp$. The canonical map $e(p) \eqdef (\downarrow p)^{\perp\perp}$ embeds $\mathbb{P}$ densely into $\mathbb{B} \setminus \{0_{\mathbb{B}}\}$ while preserving order and incompatibility ($p \perp q \iff e(p) \wedge e(q) = 0_{\mathbb{B}}$).

In both \texttt{MKTD1-theory} and \texttt{MKTD2-theory}, a Dunaev-tagged $\mathbb{B}$-name $u \in V^{(\mathbb{B})}$ is a set of Dunaev Tagged Ordered Pairs $\pair{\sigma}{b}$ with $\sigma \in V^{(\mathbb{B})}$ and $b \in \mathbb{B}$ (rejecting Kuratowski pairs). The Boolean truth values $\llbracket u \in v \rrbracket_{\mathbb{B}}$ and $\llbracket u = v \rrbracket_{\mathbb{B}}$ are defined by mutual well-founded recursion:
\[
\llbracket u \in v \rrbracket_{\mathbb{B}} \eqdef \bigvee_{\pair{\sigma}{b} \in v} \big(b \wedge \llbracket u = \sigma \rrbracket_{\mathbb{B}}\big), \qquad
\llbracket u = v \rrbracket_{\mathbb{B}} \eqdef \bigwedge_{\pair{\sigma}{b} \in u} \big(b \Rightarrow \llbracket \sigma \in v \rrbracket_{\mathbb{B}}\big) \wedge \bigwedge_{\pair{\tau}{c} \in v} \big(c \Rightarrow \llbracket \tau \in u \rrbracket_{\mathbb{B}}\big),
\]
under which all ZFC axioms evaluate to $1_{\mathbb{B}}$ and the forcing relation $p \Vdash \phi$ over Dunaev-tagged $\mathbb{P}$-names $V^{\mathbb{P}}$ satisfies $p \Vdash \phi \iff e(p) \le \llbracket \phi \rrbracket_{\mathrm{ro}(\mathbb{P})}$.

\ifdefined\qwertyflag
\undef\qwertyflag
\end{document}
\fi
